\section{提案手法}

% ===== アーキテクチャ図（後日差し替え）=====
% 図ができたら下の \fbox{...} を消して \includegraphics の行を有効にする．
% 1段に収めたい場合は figure* を figure に変える．
\begin{figure*}[t]
\centering
% \includegraphics[width=\linewidth]{img/overview}
\fbox{\begin{minipage}[c][42mm][c]{0.98\linewidth}\centering
  \textsf{[アーキテクチャ図をここに挿入]}
\end{minipage}}
\caption{提案手法の全体像}
\label{fig:overview}
\end{figure*}

本章では，REST API誤用を含むプログラムを修正するため，仕様の差分グラフを用いた自動修正手法を提案する．提案手法の全体像を\textbf{図\ref{fig:overview}}に示す．

\subsection{仕様差分グラフの構築}
非推奨仕様と最新仕様の公開ドキュメントを取得し，エンドポイント単位に正規化する．エンドポイント，パラメータ，ヘッダ，認証方式等をそれぞれノードとし，認証方式の要求，要素間の依存，値の生成手続きの3種類の関係をエッジとする．

次に，非推奨仕様の要素と最新仕様の要素の対応を宣言した移行表を記述する．この移行表と仕様グラフから，エンドポイント単位の差分ノードを生成する．差分ノードは，非推奨仕様側の要素へのFROMエッジと，最新仕様側の要素へのTOエッジをもつ．各差分ノードには，非推奨仕様側のシグネチャと変更内容を連結したテキストの埋め込みを付与する．

\subsection{誤用コードの修正}
まず誤用コードから，HTTPメソッド，パスパラメータ，ヘッダ，クエリパラメータ，リクエストボディのキーの組をプログラム要素として抽出する．
% 対象データセットが8言語を含むため，抽出は言語非依存の正規表現に基づく．1つのコードから複数のプログラム要素が得られる．
% →これは今後tree-sitterでAST構築したいし言わない．

% 微妙，消してもいい
% 各プログラム要素について，まず非推奨仕様側のシグネチャとの決定的なマッチを試みる．仕様がパスパラメータを\verb|[user-id]|のように記述するのに対しコードは変数展開や具体値を用いるため，セグメント単位でパスパラメータをワイルドカードとして照合する．また，コードが既に最新仕様のパスを用いている場合は，その要素に差分を適用しない．マッチしなかったプログラム要素についてのみ，差分ノードの埋め込みに対するベクトル検索へ回す．\textbf{検索対象は差分ノードのみであり，仕様ノードを直接検索することはない．}

引かれた差分ノードからは，TOエッジの先のみを探索し，最新仕様で必要となる仕様要素を収集する．SwitchBot APIの例では，v1.1のエンドポイントから認証方式ノードへ辿り，そこから\texttt{t}，\texttt{sign}，\texttt{nonce}の各ヘッダノードと，HMAC-SHA256による署名生成の手続きノードが収集される．収集結果はプログラム要素ごとに整形する．
% 差分カードは，変更の要約，非推奨仕様と最新仕様を対比する表，最新仕様で必要となる要素の一覧から成る．
% 同一の差分ノードが複数のプログラム要素に対応する場合は，本文を一度だけ展開し，二度目以降は参照に置き換える．
